Әлемде әр адамның өз жолы бар. Іні мен әпке бірге ертегі оқыды, ал сыртта ғажайып түн болды. Үміт пен қуаныш жүрегінде мәңгі сақталады. Ұлы дала әні шалқып, қыранның қанаты аспанды жайлады. Өмір өзені асықпай ағады, әр толқыны жаңа оқиға әкеледі. Һәм барлығы, ақыр соңында, тыныштықпен ұйықтады.

Қазақ тілінің тоғыз ерекше әрпі — Ә, І, Ң, Ғ, Ү, Ұ, Қ, Ө, Һ — осы мәтінде әдейі жиі қолданылды: диакритикалық белгілердің тураланымы, әріп биіктігі және сызық қалыңдығы негізгі қаріппен үйлесімді екенін көзбен тексеру үшін.

Таңның атысы, күннің батысы — бәрі де кезегімен келеді. Дала кеңістігінде жел есіп, шөп басы бір ырғақпен тербеледі. Осындай сәттерде адам баласы өз болмысын, өз тілін, өз топырағын тереңірек түсінеді.
